\section{ISAC-Assisted Underserved-Region Model}
\label{sec:mathematical-model}

Consider ground UEs $\mathcal{U}$, ISAC sensing nodes $\mathcal{S}$,
and UAV gNBs $\mathcal{Q}$. Communication KPIs determine which UEs
are underserved, while ISAC observations estimate their spatial
distribution. The near-RT RIC fuses both sources, predicts future
underserved regions, and supplies target centroids to the UAV
repositioning controller.

\subsection{Sensing and Communication Fusion}

For sensing node $s$ and UE $u$, the sensing range and system-level
echo SNR are
\begin{align}
 r_{s,u}(t)&=\|\mathbf{p}_s(t)-\mathbf{p}_u(t)\|_2, \\
 \gamma_{s,u}(t)&=
 \frac{P_sG_s^2\lambda^2\sigma_u}
 {(4\pi)^3r_{s,u}^4(t)L_sN_0B_s},
 \label{eq:isac-snr-compact}
\end{align}
where $P_s$, $G_s$, $\lambda$, $\sigma_u$, $L_s$, $N_0$, and
$B_s$ denote sensing power, antenna gain, wavelength, radar cross
section, aggregate loss, noise density, and sensing bandwidth,
respectively. Detection is sampled according to
\begin{equation}
 d_{s,u}(t)\sim\operatorname{Bernoulli}\!\left(
 \frac{1}{1+\exp[-a(\gamma_{s,u}^{\mathrm{dB}}(t)-\gamma_0)]}
 \right).
 \label{eq:isac-detection-compact}
\end{equation}
For $d_{s,u}(t)=1$, the sensed position is
\begin{equation}
 \widehat{\mathbf{p}}_{s,u}(t)=\mathbf{p}_u(t)+\mathbf{e}_{s,u}(t),
 \quad
 \mathbf{e}_{s,u}(t)\sim\mathcal{N}
 (\mathbf{0},\sigma_{s,u}^2(t)\mathbf{I}),
 \label{eq:isac-position-compact}
\end{equation}
where
\begin{equation}
 \sigma_{s,u}(t)=\operatorname{clip}\!\left[
 \sigma_{\mathrm{ref}}
 10^{-\frac{\gamma_{s,u}^{\mathrm{dB}}(t)-\gamma_{\mathrm{ref}}}{20}},
 \sigma_{\min},\sigma_{\max}\right].
 \label{eq:isac-error-compact}
\end{equation}
Multiple detections are fused at the RIC using
\begin{equation}
 \widehat{\mathbf{p}}_u(t)=
 \frac{\sum_{s\in\mathcal{S}_u(t)}
 \sigma_{s,u}^{-2}(t)\widehat{\mathbf{p}}_{s,u}(t)}
 {\sum_{s\in\mathcal{S}_u(t)}\sigma_{s,u}^{-2}(t)},
 \label{eq:isac-fusion-compact}
\end{equation}
where $\mathcal{S}_u(t)=\{s:d_{s,u}(t)=1\}$. True positions are
retained only for simulation evaluation.

Let $R_u$, $P_u$, $C_u$, and $D_u$ denote serving RSRP, PDR,
throughput, and delay. The communication-derived service indicator is
\begin{equation}
 I_u(t)=\mathbf{1}\!\left\{
 R_u<R_{\min}\lor P_u<P_{\min}\lor
 C_u<C_{\min}\lor D_u>D_{\max}\right\}.
 \label{eq:underserved-compact}
\end{equation}
Only KPIs measured in the experiment are applied.

\subsection{Underserved Heatmap and RF Prediction}

For grid cell $g\in\mathcal{G}$ with centre $\mathbf{c}_g$, sensed
density, underserved intensity, and underserved share are
\begin{align}
 \rho_g(t)&=\sum_{u}K_h(\widehat{\mathbf{p}}_u^{xy}-\mathbf{c}_g),\\
 H_g(t)&=\sum_{u}I_u(t)
 K_h(\widehat{\mathbf{p}}_u^{xy}-\mathbf{c}_g),\\
 Z_g(t)&=\frac{H_g(t)}{\rho_g(t)+\epsilon}.
 \label{eq:heatmap-compact}
\end{align}
The RF uses the available subset of
\begin{equation}
 \mathbf{x}_g(t)=[\rho_g,H_g,Z_g,\overline{R}_g,R_g^{\min},
 \overline{P}_g,\overline{C}_g,\overline{D}_g,
 \overline{L}_g,\overline{v}_g,\overline{\sigma}_g]^T.
 \label{eq:rf-features-compact}
\end{equation}
For prediction horizon $\Delta_p$, oracle quantities
$N_g^\star$ and $Z_g^\star$ are used only to create the training label
\begin{equation}
 y_g(t)=\mathbf{1}\!\left\{
 N_g^\star(t+\Delta_p)\geq N_{\min}\land
 Z_g^\star(t+\Delta_p)\geq\beta\right\}.
 \label{eq:rf-label-compact}
\end{equation}
The RF estimates $q_g(t)=\Pr[y_g(t)=1\mid\mathbf{x}_g(t)]$, and
$\widehat{\mathcal{G}}_{\mathrm{US}}(t)=
\{g:q_g(t)\geq\tau\}$. Training, validation, and testing use
independent simulation seeds to prevent temporal leakage.

\subsection{Clustering and UAV Repositioning}

For $K\leq|\mathcal{Q}|$ available UAVs, probability-weighted
K-means obtains target centroids:
\begin{equation}
 \min_{\{\boldsymbol{\mu}_k\}_{k=1}^{K}}
 \sum_{g\in\widehat{\mathcal{G}}_{\mathrm{US}}}
 q_gH_g\min_k\|\mathbf{c}_g-\boldsymbol{\mu}_k\|_2^2.
 \label{eq:kmeans-compact}
\end{equation}
Each UAV is assigned once to an unassigned centroid by minimum travel
distance. For assigned centroid $\boldsymbol{\mu}_{k(j)}$, UAV $j$
updates its horizontal position as
\begin{equation}
 \mathbf{q}_j(t+\Delta t)=\mathbf{q}_j(t)+
 \min\{v_{\max}\Delta t,\|\boldsymbol{\delta}_j\|_2\}
 \frac{\boldsymbol{\delta}_j}{\|\boldsymbol{\delta}_j\|_2},
 \quad
 \boldsymbol{\delta}_j=\boldsymbol{\mu}_{k(j)}-\mathbf{q}_j(t),
 \label{eq:movement-compact}
\end{equation}
with zero displacement when $\|\boldsymbol{\delta}_j\|_2=0$ and
subject to flight-area and altitude constraints.

Sensing performance is evaluated using detection probability,
localisation RMSE, hotspot F1, and matched centroid error. The
end-to-end benefit is evaluated using the underserved-UE fraction,
coverage probability, fifth-percentile throughput, PDR, and outage
time before and after repositioning.
